% !PRINTED 402
\weberpage{402}

Auch ist hier noch darauf hinzuweisen, dass aus der Realit\"at
einer Doppeltangente noch keineswegs die Realit\"at der Ber\"uhrungspunkte
 folgt, weil diese Ber\"uhrungspunkte wieder von einer
quadratischen Gleichung abh\"angen. Die reellen Doppeltangenten
zerfallen also wieder in zwei Arten, solche mit reellen und
solche mit imagin\"aren Ber\"uhrungspunkten. Was hier f\"ur m\"ogliche
 F\"alle zu unterscheiden sind, diese Frage er\"ortern wir nicht
weiter \textsuperscript{1}).

\begin{center}
\rule{0.3\linewidth}{0.4pt}
\end{center}

\textsuperscript{1}) Vergl. \"uber die ganze Frage von der Realit\"at der Doppeltangenten
vom geometrischen Gesichtspunkte: \gesperrt{Zeuthen}, \glqq Sur les diff\'erentes formes
des courbes planes du quatri\`eme ordre\grqq{}. Mathem. Ann., Bd. VII (1873).
